\documentclass[a4paper,12pt]{article}
\usepackage[utf8]{inputenc}
\usepackage[T1]{fontenc}
\usepackage[ngerman]{babel}
\usepackage{geometry}
\geometry{margin=2.5cm}
\begin{document}

\section*{Applying Physics-Informed Neural Networks to Solve Navier--Stokes
Equations for Laminar Flow around a Particle}

\textbf{Autoren:} Beichao Hu, Dwayne McDaniel \\
\textbf{Jahr:} 2023 \\
\textbf{Zeitschrift:} Mathematical and Computational Applications, 28(5), 102 \\
\textbf{DOI:} \texttt{https://doi.org/10.3390/mca28050102}

% ------------------------------------------------------------------
\subsection*{Zusammenfassung}
% ------------------------------------------------------------------

Diese Arbeit entwickelt ein Physics-Informed Neural Network (PINN) zur Lösung
der zweidimensionalen, stationären, inkompressiblen Navier-Stokes-Gleichungen
für laminare Umströmungen einfacher Partikelgeometrien. Als Benchmark dient
die Umströmung eines kreisförmigen Partikels; zusätzlich wird ein elliptischer
Partikel untersucht, um die Robustheit des Ansatzes gegenüber veränderten
Geometrien zu evaluieren. Die Simulationen werden bei Reynolds-Zahlen von
$Re = 5$, $20$ und $50$ durchgeführt, also im laminaren Regime, in dem
Viskositätskräfte und Trägheitskräfte von vergleichbarer Größenordnung sind.

Das verwendete Netzwerk ist ein vollständig verbundenes (Fully Connected) PINN
mit zwei Eingangsneuronen ($x$- und $y$-Koordinate), acht versteckten Schichten
mit je 60 Neuronen und drei Ausgangsneuronen ($u$, $v$, $p$). Die physikalischen
Gesetze -- Massenerhaltung sowie Impulserhaltung in $x$- und $y$-Richtung --
werden direkt in die Verlustfunktion eingebettet, sodass kein beschriftetes
Trainingsdaten\-set erforderlich ist. Randbedingungen (Einströmrand, Ausströmrand,
Wände) werden als zusätzliche Terme in der Verlustfunktion weich erzwungen und
mit Strafgewichten $\alpha$ versehen, um die Balance zwischen den einzelnen
Verlustbeiträgen zu steuern. Das Training erfolgt zunächst mit dem Adam-Optimizer
(10\,000 Iterationen) und anschließend mit L-BFGS zur Feinoptimierung; die
Konvergenz wird durch einen Gradienten-Schwellenwert von $10^{-7}$ definiert.

Die Ergebnisse zeigen, dass das PINN die wesentlichen Strömungsmerkmale --
Staupunkt, Venturi-Effekt, Nachlaufgebiet und Grenzschicht -- bei allen
untersuchten Reynolds-Zahlen qualitativ korrekt reproduziert. Quantitativ
liegt der Fehler im berechneten Widerstandsbeiwert $C_d$ in allen Szenarien
unter 10\,\% gegenüber CFD-Referenzlösungen (ANSYS Fluent 2020R2). Der
gesamte Widerstandsbeiwert ergibt sich aus einem Druckanteil und einem
viskosen Anteil, die jeweils über Oberflächenintegrale bestimmt werden, da
das punktbasierte PINN-Verfahren keine Gittervolumina bereitstellt. Mit
zunehmender Reynolds-Zahl steigen die Abweichungen leicht an, was die Autoren
auf das weiche Erzwingen der Randbedingungen und eine dadurch bedingte
unvollständige Massenerhaltung zurückführen.

Ein praktisch bedeutsames Ergebnis betrifft die Rechenzeit: Das Training
benötigt auf einer Nvidia Tesla P100 GPU etwa 1{,}5 Stunden, während die
Inferenz für das gesamte Berechnungsgebiet in 0{,}07 Sekunden abgeschlossen
ist. CFD (CPU) benötigt für denselben Fall 2 Minuten. Damit ist das PINN
in der Trainingsphase deutlich teurer, bietet aber in der Anwendungsphase
einen erheblichen Geschwindigkeitsvorteil.

% ------------------------------------------------------------------
\subsection*{Methodische Kernpunkte}
% ------------------------------------------------------------------

\begin{itemize}
    \item \textbf{Verlustfunktion:} Summe aus PDE-Residuen (Kontinuität,
    $x$- und $y$-Impuls) und Randbedingungsresiduen; Strafgewichte $\alpha$
    kompensieren die ungleiche Verteilung der Trainingspunkte.
    \item \textbf{Automatische Differentiation (AD):} Ableitungen werden
    exakt auf dem Berechnungsgraphen berechnet, ohne Diskretisierungsfehler
    wie bei klassischen FVM- oder FDM-Verfahren.
    \item \textbf{Widerstandsbeiwert:} Berechnung durch numerische Integration
    von Druckfeld und viskosem Spannungstensor entlang der Partikeloberfläche
    (Trapezregel).
    \item \textbf{Trainingsstrategie:} Adam + L-BFGS in Kombination; L-BFGS
    erhöht die Konvergenzstabilität in der Feinoptimierungsphase erheblich.
\end{itemize}

% ------------------------------------------------------------------
\subsection*{Kritische Einschätzung}
% ------------------------------------------------------------------

Die Arbeit demonstriert die grundsätzliche Eignung von PINNs für die
Partikelumströmung bei niedrigen Reynolds-Zahlen, weist aber auch auf
inhärente Schwächen des Ansatzes hin. Das weiche Erzwingen der
Randbedingungen führt zu leichten Massenerhaltungsfehlern, die mit
steigender Reynolds-Zahl zunehmen. Die Trainingszeit ist im Vergleich zu
klassischer CFD deutlich höher, was die Einsatzfähigkeit als schnellen
Ersatzmodel einschränkt -- insbesondere wenn das Netzwerk für jede neue
Geometrie erneut trainiert werden muss. Darüber hinaus beschränkt sich die
Arbeit auf einen einzelnen Partikel in einem kanalförmigen Gebiet; eine
Generalisierung auf variable Partikelanzahlen oder -positionen wird nicht
untersucht. Die Autoren selbst benennen dies als offenen Forschungsbedarf
für zukünftige Arbeiten.

% ------------------------------------------------------------------
\subsection*{Bezug zur Masterarbeit}
% ------------------------------------------------------------------

Diese Arbeit ist für die vorliegende Masterarbeit aus mehreren Gründen
relevant, unterscheidet sich aber in grundlegenden Punkten von deren Ansatz.

\textbf{Gemeinsamkeiten:} Das physikalische Setting ist eng verwandt. Die
Umströmung eines einzelnen Partikels bei kleinen Reynolds-Zahlen entspricht
dem Stokes-Strömungsregime, das für die Kristallsedimentation in Magmakammern
charakteristisch ist ($Re \ll 1$). Beide Arbeiten zielen darauf ab, den
rechenintensiven Einsatz klassischer CFD-Verfahren durch datengestützte
Modelle zu beschleunigen, und nutzen physikalisch motivierte Verlustterme
zur Verbesserung der Lösungsqualität. Die in dieser Arbeit diskutierte
Herausforderung, Massenerhaltung durch weiches Erzwingen der
Randbedingungen zu gewährleisten, ist direkt mit der Verwendung der
Stromfunktions-Formulierung in der Masterarbeit verknüpft: Durch die
Vorhersage der Stromfunktion $\psi$ anstelle der Geschwindigkeitskomponenten
wird Inkompressibilität in der Masterarbeit mathematisch garantiert, was die
in Hu \& McDaniel beschriebenen Massenerhaltungsfehler grundsätzlich vermeidet.

\textbf{Wesentliche Unterschiede:} Während Hu \& McDaniel ein einzelnes
PINN verwenden, das für eine fixe Geometrie trainiert wird und keine
Trainingsdaten benötigt, verfolgt die Masterarbeit einen datenbasierten
Surrogatmodell-Ansatz (U-Net, FNO), der auf LaMEM-Simulationen als
Grundwahrheit trainiert wird. Der zentrale Beitrag der Masterarbeit --
die Generalisierung über variable Kristallanzahlen (1 bis $N$) und
räumliche Anordnungen -- ist in Hu \& McDaniel nicht adressiert. Das PINN
in dieser Arbeit ist auf eine Instanz beschränkt und müsste für jede neue
Konfiguration von Grund auf neu trainiert werden, was es für den angestrebten
Einsatz in geowissenschaftlichen Parameterstudien ungeeignet macht. Die
gitterbasierten Architekturen der Masterarbeit (U-Net, FNO auf einem
$256 \times 256$-Gitter) ermöglichen demgegenüber eine einmalige
Trainingsphase mit anschließender instantaner Inferenz für beliebige
Kristallkonfigurationen innerhalb des Trainingsbereichs.

\textbf{Einordnung in den Literaturkontext:} Die vorliegende Arbeit belegt,
dass PINNs für die Umströmung einfacher Partikelgeometrien prinzipiell
geeignet sind, verdeutlicht aber zugleich deren Limitierungen hinsichtlich
Trainingsaufwand und Generalisierungsfähigkeit. Dies motiviert den in der
Masterarbeit gewählten Ansatz, auf diskrete Simula\-tionsdaten und
skalierbare neuronale Operatoren zu setzen, um die Anforderungen variabler
Mehr\-partikelkonfigurationen zu erfüllen.

\end{document}